\section{Data Preprocessing and Exploration}\label{sec:data}

We use only the official files in \texttt{ProblemC\_Music\_data/}; no
external data are introduced. Of the four files,
\texttt{influence\_data.csv} ($42{,}770$ directed influencer$\to$follower
edges, with both artists' \texttt{main\_genre} and \texttt{active\_start}
decade) is the primary input to Model~I and also serves as the genre
dictionary for Model~II, and \texttt{data\_by\_artist.csv} ($5{,}810$
artists) provides artist-level mean Spotify features for Model~II; the
song-level \texttt{full\_music\_data.csv} and yearly
\texttt{data\_by\_year.csv} are reserved for finer temporal analyses and
are not used in the models below. This section documents the two cleaning
pipelines and two data-quality audits whose findings shape every later
model; the influence data require their own cleaning, and the feature
data a different one, so we keep the two pipelines explicit and separate.

\subsection{Cleaning of the Influence Data}\label{subsec:data-cleaning}

The raw influence file requires four repairs before it can be used.
(i)~\emph{Identifier hygiene}: trailing whitespace in artist identifiers and
names is trimmed. (ii)~\emph{Encoding artefacts}: mojibake in some artist
names is repaired by re-decoding the affected strings. (iii)~\emph{Label
repair}: genre labels contain an HTML-escape residue (e.g.\ ``R\&B;''
$\to$ ``R\&B''), which we strip. (iv)~\emph{Type coercion}: active-start
fields are cast to decade integers. After these repairs we check for
duplicated directed pairs; none exist, so all 42{,}770 unique directed pairs
are retained, spanning $N=5{,}603$ artists.

\subsection{Cleaning Choices for Model~II}\label{subsec:data-cleaning-features}

\begin{enumerate}
  \item Join \texttt{data\_by\_artist} to genre labels via
        \texttt{artist\_id}, where genres are taken from
        \texttt{influence\_data} (we do \emph{not} redefine Pop/Rock or
        other genres; see Assumption~A4).
  \item Drop artists with missing genre or genre \texttt{Unknown}.
  \item Select $p=11$ numeric features: danceability, energy, valence,
        tempo, loudness, acousticness, instrumentalness, liveness,
        speechiness, duration\_ms, popularity.
  \item Apply column-wise $z$-score standardization before computing
        Euclidean distances (see Model~II). We note that
        \emph{popularity} is not comparable across eras---it rises almost
        monotonically from 1921 to 2020 in \texttt{data\_by\_year}---so a
        within-year normalization variant is examined as a robustness check
        in Model~II.
\end{enumerate}
After cleaning we retain $n=5{,}599$ artists with both features and a genre
label, covering 19 genres (Pop/Rock is the largest class). This similarity
sample ($5{,}599$ artists) is not identical to the network sample
($5{,}603$ artists of \cref{subsec:data-cleaning}); the small discrepancy
arises because having a genre label, having Spotify features, and appearing
in the influence data are three separate conditions, and any analysis that
crosses the two samples uses their intersection (stated explicitly where it
occurs).

\subsection{Data-Quality Audit}\label{subsec:data-audit}

Two structural facts of the influence data shape every later model, so we
audit them here rather than inside any single model.

\emph{Time-order consistency.} Decade information allows a consistency
check that the raw data never advertises: in 733 edges (1.7\%) the
follower's active-start decade precedes the influencer's, and in 8{,}477
edges the two decades coincide. Per Assumption~A3 the 733 reversals are
retained in the primary network and their effect is quantified in
\cref{subsec:q1-robustness}; same-decade edges are kept as same-generation
mutual influence.

\emph{Skewed connectivity and genre concentration.} The out-degree
distribution is heavily right-skewed (mean 11.3, median 5, maximum 615 for
The Beatles), while in-degrees are moderate (median 7, maximum 47): being
cited is bounded by list length, citing is not. Genre labels concentrate
sharply---Pop/Rock accounts for 56\% of all edges---and 76.6\% of edges
connect same-genre artists. These facts motivate the log-scale treatment of
degree in \cref{subsec:q1-global}, the cohort view of
\cref{subsec:q1-cohorts}, and the genre-comparison design of Model~III.
